% Vietnamese translation for OI Public Library
% Source-PDF: 比赛 CONTEST/2021“MINIEYE杯”中国大学生算法设计超级联赛/2021“MINIEYE杯”中国大学生算法设计超级联赛（6）/2021“MINIEYE杯”中国大学生算法设计超级联赛（6）-题解集.pdf
\documentclass[11pt]{article}
\usepackage{oipl-vn}
\usepackage{tikz}

\title{Lời giải trận thứ sáu HDU Multi-University 2021}
\author{}
\date{}

\begin{document}
\maketitle

\origtitle{HDU Multi-University 2021 MINIEYE Cup, Contest 6 Solutions}
\sourcepath{contest/hdu-2021-minieye-06-solutions.pdf}

\section{1001 Yes, Prime Minister}

Rõ ràng $r>0$ và $r\ge l$. Nếu $l\le 0$, tổng trên đoạn $[l,r]$ bằng tổng
trên đoạn $[-l+1,r]$, đồng thời $r\ge -l+1>0$. Vì vậy với mọi đoạn, luôn có
thể tìm một đoạn tương ứng thỏa $r\ge l>0$ và có cùng tổng.

Xét
\[
  \sum_{i=l}^{r} i
  = \frac{(l+r)(r-l+1)}{2}.
\]
Khi $l>0$, nếu $r-l\ge 2$ thì trong hai số
\[
  \frac{l+r}{2},\qquad \frac{r-l+1}{2}
\]
chắc chắn có một số nguyên lớn hơn $1$. Do đó tổng đoạn tách được thành tích
của hai thừa số, nên không thể là số nguyên tố. Vì vậy mọi đoạn có tổng là
số nguyên tố đều có độ dài không quá $2$.

Tiền xử lý các số nguyên tố trong một khoảng hơi lớn hơn $[2,2\cdot 10^7]$;
có thể dùng sàng Eratosthenes hoặc sàng tuyến tính.

Nếu $x\le 0$, các đáp án ứng viên là:
\begin{itemize}
  \item số nhỏ nhất $y\ge 1-x$ sao cho $y$ là số nguyên tố; đoạn là
  $[-y+1,y]$;
  \item số nhỏ nhất $z\ge 2-x$ sao cho $2z-1$ là số nguyên tố; đoạn là
  $[-z+2,z]$.
\end{itemize}

Nếu $x>0$, các đáp án ứng viên là:
\begin{itemize}
  \item $[x,x]$;
  \item $[x,x+1]$;
  \item $[x-1,x]$;
  \item số nhỏ nhất $y\ge x$ sao cho $y$ là số nguyên tố; đoạn là
  $[-y+1,y]$;
  \item số nhỏ nhất $z\ge x$ sao cho $2z-1$ là số nguyên tố; đoạn là
  $[-z+2,z]$.
\end{itemize}

Sau đó chỉ cần tìm kiếm nhị phân. Độ phức tạp thời gian là
\[
  O(T\log |x|+|x|).
\]

\section{1002 Might and Magic}

Đáp án bằng lượng sát thương lớn nhất có thể gây ra khi máu đối thủ được xem
là vô hạn. Chú ý số vòng sống sót là
\[
  t=\left\lceil \frac{\mathrm{Damage}}{H_0}\right\rceil .
\]
Ta có thể dùng chia khối theo thương để liệt kê $t$ và lượng phòng thủ $D_0$
cần bỏ ra trong độ phức tạp $O(\sqrt{H_0})$.

Tiếp theo dùng nhận xét: hoặc cộng toàn bộ vào $A_0$, hoặc cộng toàn bộ vào
$P_0,K_0$. Chứng minh như sau.

Cố định số lần đánh vật lý và đánh phép là $a,b$ với $a+b=t$. Nếu $K_0$
không đủ thì xem như bỏ qua lượt. Giả sử cộng $x$ điểm vào vật lý/phép, hàm
sát thương vật lý là
\[
  P(x)=
  \begin{cases}
    aC_p, & x\le D_1,\\
    aC_p(x-D_1), & x>D_1,
  \end{cases}
\]
là một hàm lồi. Hàm sát thương phép là
\[
  M(x)=
  \begin{cases}
    C_m\left(x-\left\lfloor\frac{x}{2}\right\rfloor\right)
       \left\lfloor\frac{x}{2}\right\rfloor, & 2x\le b,\\
    C_m(x-b)b, & 2x>b,
  \end{cases}
\]
cũng là một hàm lồi. Vì thế tổng sát thương
\[
  F(x)=P(x)+M(N-D_0-x)
\]
theo số điểm cộng vào vật lý $x$ cũng là hàm lồi, nên giá trị lớn nhất
chắc chắn đạt được tại ít nhất một trong hai đầu mút.

Trong trường hợp cộng toàn bộ vào vật lý, chỉ cần cộng hết vào $A_0$.
Trong trường hợp cộng toàn bộ vào phép, cần tìm giá trị lớn nhất của
\[
  D(K_0)=C_mK_0(N-D_0-K_0)+C_p(t-K_0),
  \qquad 0\le K_0\le \min(t,N-D_0).
\]

Độ phức tạp thời gian là $O(T\sqrt{H_0})$.

\section{1003 0 tree}

Xét việc chuyển toàn bộ trọng số cạnh thành trọng số điểm. Đặt trọng số điểm
dạng tổng là
\[
  b_i=\sum_{e=\langle i,v\rangle} b_e.
\]
Gọi $a_i$ là trọng số điểm dạng xor, và $b_i$ là trọng số điểm dạng tổng. Có
thể thấy việc đưa toàn bộ trọng số cạnh của đồ thị về $0$ tương đương với
đưa cả hai loại trọng số điểm về $0$.

Khi thực hiện thao tác $x\ y\ w$, ngoài hai đầu mút $x,y$, mọi trọng số điểm
trên đường đi đều không đổi. Nếu $\operatorname{dis}(x,y)$ chẵn, thao tác
tương đương với
\[
  a_x\gets a_x\oplus w,\qquad
  a_y\gets a_y\oplus w,\qquad
  b_x\gets b_x+w,\qquad
  b_y\gets b_y-w.
\]
Nếu $\operatorname{dis}(x,y)$ lẻ, thao tác tương đương với
\[
  a_x\gets a_x\oplus w,\qquad
  a_y\gets a_y\oplus w,\qquad
  b_x\gets b_x+w,\qquad
  b_y\gets b_y+w.
\]

Trước hết tô màu $01$ cho cây. Không có nghiệm nếu thỏa bất kỳ điều kiện nào
sau đây:
\begin{itemize}
  \item Điều kiện xor theo hai màu không khớp:
  \[
    \bigoplus_{\operatorname{col}_i=0} a_i
    \ne
    \bigoplus_{\operatorname{col}_i=1} a_i.
  \]
  \item $a_i$ và $b_i$ khác tính chẵn lẻ. Lý do là trong suốt quá trình thao
  tác, tính chẵn lẻ của $a_i$ và $b_i$ thay đổi đồng bộ.
  \item Điều kiện tổng trọng số theo màu không đủ để bù phần xor:
  \[
    \bigoplus_{\operatorname{col}_i=k} a_i
    >
    -\sum_{\operatorname{col}_i=k} b_i.
  \]
  Các thao tác giữa hai điểm cùng màu không đổi tổng trọng số điểm, nên thao
  tác giữa hai màu khác nhau phải có khả năng đưa cả hai tổng trọng số về
  $0$. Vì tổng của một số số luôn không nhỏ hơn xor của chúng, ta nhận được
  điều kiện trên.
\end{itemize}

Trước hết đưa các trọng số dạng tổng về $0$. Gọi xor của các trọng số xor
trong cùng một màu là $A$, và tổng các trọng số dạng tổng là $B$; khi đó
$A\le -B$. Có thể chọn hai điểm khác màu rồi lần lượt thực hiện các thao tác
với
\[
  w=A,\qquad w=-\frac{A+B}{2},\qquad w=-\frac{A+B}{2}
\]
để đạt được mục tiêu này.

Với hai điểm cùng màu bất kỳ $x,y$, ta có thể thực hiện hai lần thao tác
$x\ y\ v$. Hiệu quả là không đổi $a_i$, còn với hai điểm cùng màu:
\[
  b[x]\mathrel{+}=2v,\qquad b[y]\mathrel{-}=2v.
\]
Ta cũng có thể thực hiện hai lần thao tác $x\ y\ v$, rồi thêm một lần thao
tác $y\ x\ 2v$. Hiệu quả là không đổi $b_i$, còn với hai điểm cùng màu:
\[
  a[x]\gets a[x]\oplus 2v,\qquad
  a[y]\gets a[y]\oplus 2v.
\]

Vì vậy trước hết đưa thô một trong hai loại trọng số điểm về $0$, sau đó dùng
một trong hai phương pháp trên để đưa loại còn lại về $0$. Số thao tác theo
cách này không vượt quá $3n$ hoặc $4n$. Độ phức tạp thời gian và bộ nhớ là
$O(n)$.

\section{1004 Decomposition}

Xét cách dựng một chu trình đơn như hình minh họa dưới đây. Cách dựng là xuất
phát từ điểm trung tâm, chia các điểm còn lại thành hai nhóm trái trên và
phải dưới, nối xen kẽ điểm hiện tại với các điểm thuộc hai nhóm, rồi cuối
cùng quay lại điểm trung tâm.

\begin{center}
\begin{tikzpicture}[scale=.72,line cap=round,line join=round]
  \tikzset{
    outer/.style={circle,draw,fill=gray!8,minimum size=5.5mm,inner sep=0pt},
    middle/.style={circle,draw,fill=blue!12,minimum size=5.5mm,inner sep=0pt}
  }
  \foreach \x/\ang in {0/0,5.8/120,11.6/240} {
    \begin{scope}[shift={(\x,0)},rotate=\ang]
      \draw[very thick,magenta!85!black]
        (0,0) -- (0,2.15) -- (1.85,1.05) -- (-1.85,1.05) --
        (1.85,-1.05) -- (-1.85,-1.05) -- (0,-2.15) -- cycle;
      \node[middle] at (0,0) {};
      \foreach \p in {(0,2.15),(1.85,1.05),(-1.85,1.05),
                     (1.85,-1.05),(-1.85,-1.05),(0,-2.15)}
        \node[outer] at \p {};
    \end{scope}
  }
  \draw[blue!70!black,densely dashed,very thick] (2.9,-2.8) -- (2.9,2.8);
  \draw[blue!70!black,densely dashed,very thick] (8.7,-2.8) -- (8.7,2.8);
  \node at (0,-3.25) {cấu hình cơ sở};
  \node at (5.8,-3.25) {sau một lần quay};
  \node at (11.6,-3.25) {sau hai lần quay};
\end{tikzpicture}
\end{center}

Có thể chứng minh rằng bằng cách liên tục quay cấu hình này một góc
\[
  2\pi\cdot\frac{2}{n-1}
\]
ta thu được $(n-1)/2$ chu trình khác nhau, và các chu trình đó đi qua đúng
một lần mỗi cạnh của đồ thị đầy đủ. Ghép các chu trình theo thứ tự quay sẽ
thu được một chu trình Euler của đồ thị đầy đủ. Có thể chứng minh rằng mọi
$n-2$ điểm liên tiếp trong chu trình Euler này đều đôi một khác nhau, nên ta
chỉ cần cắt trực tiếp chu trình Euler đó thành các đường đi theo thứ tự đề
bài yêu cầu.

Ý tưởng chứng minh tính chất thứ nhất: trước hết, các cạnh giữa điểm trung
tâm và những điểm khác hiển nhiên đều đã được thăm. Tiếp theo, nếu đánh số
các điểm còn lại theo chiều kim đồng hồ như các phần tử modulo $n-1$, thì
trong mỗi chu trình đơn, hiệu số chỉ số của hai đầu mút trên các cạnh không
đi qua điểm trung tâm lần lượt là
\[
  +1,-2,+3,-4,\ldots,+(n-2).
\]
Theo nghĩa đồng dư, đây đúng là
\[
  \pm 1,\pm 2,\ldots,\pm\frac{n-1}{2}.
\]
Phân tích thêm cho thấy khi quay, điểm bắt đầu của mỗi loại trị tuyệt đối
của hiệu số sẽ lần lượt nhận mọi giá trị $0,1,\ldots,n-2$. Do đó cách dựng
không lặp cạnh và đi qua mỗi cạnh của đồ thị đầy đủ đúng một lần.

Ý tưởng chứng minh tính chất thứ hai: điểm trung tâm xuất hiện hai lần liên
tiếp cách nhau đúng $n$ vị trí; với các điểm khác, phân tích tương tự hình
minh họa cho thấy khoảng cách giữa hai lần xuất hiện không nhỏ hơn $n-1$.

Nhận xét: ý tưởng này cũng có thể dùng để dựng cách phân rã đồ thị đầy đủ
thành các đường Hamilton, các ghép cặp hoàn hảo, v.v.

\section{1005 Median}

Rõ ràng $m$ số $b_1,b_2,\ldots,b_m$ phải được đặt vào $m$ tập khác nhau. Các
số còn lại, tức $n-m$ số, cần được đặt vào $m$ tập này sao cho không ảnh
hưởng đến trung vị của mỗi tập.

Dùng một ví dụ để dễ giải thích: giả sử $n=6$, $m=2$, $b_1=3$, $b_2=5$. Khi
đó các số $1,2,\ldots,n$ bị dãy $b$ chia thành ba đoạn $1,2$, $4$, và $6$;
bất kỳ hai số thuộc hai đoạn khác nhau có thể ghép cặp và triệt tiêu lẫn
nhau. Vì vậy toàn bộ số còn sót lại cuối cùng nhất định thuộc cùng một đoạn.
Ta xét hai trường hợp.

\begin{itemize}
  \item Nếu độ dài đoạn lớn nhất không vượt quá tổng độ dài các đoạn còn
  lại, cuối cùng hoặc mọi số đều bị triệt tiêu, hoặc còn lại đúng một số, và
  số còn lại này có thể nằm trong bất kỳ đoạn nào. Nếu có số còn lại, có thể
  giả sử để lại một số ở đoạn cuối; khi đó đặt các số của đoạn cuối vào tập
  có trung vị nhỏ nhất là đủ thỏa yêu cầu. Vì thế đáp án là
  \texttt{YES}.
  \item Nếu độ dài đoạn lớn nhất lớn hơn tổng độ dài các đoạn còn lại, toàn
  bộ số còn lại cuối cùng đều nằm trong đoạn lớn nhất. Gọi $x$ là số tập có
  trung vị nhỏ hơn giá trị nhỏ nhất của đoạn này. Dễ thấy có nghiệm khi và
  chỉ khi $x$ không nhỏ hơn số lượng số còn lại trong đoạn đó.
\end{itemize}

Độ phức tạp thời gian là $O(n)$.

\section{1006 The Struggle}

Độ phức tạp thuật toán của bài này là $O(n\log n)$, với
\[
  n=\max_{(x,y)\in E}\max(x,y).
\]
Xét cách tính bằng tích chập theo hiệu đối xứng của tập. Mỗi lần tích chập
theo hiệu đối xứng chỉ xử lý được một khối dạng
\[
  [x\cdot 2^k,\,x\cdot 2^k+2^k-1]
  \times
  [y\cdot 2^k,\,y\cdot 2^k+2^k-1].
\]
Ta xử lý trước mọi khối lớn nhất nằm hoàn toàn trong ellipse, rồi đến các
khối lớn thứ hai, và cứ tiếp tục như vậy.

Thuật toán trực tiếp theo mô tả trên có hai thừa số log, vẫn chưa đủ nhanh.
Cách tối ưu là thực hiện FWT từ dưới lên; khi đang ở tầng nào thì xử lý các
khối cần tính của tầng đó. Sau khi tính được tích, không cần vội thực hiện
FWT ngược, mà cộng dồn luôn vào mảng kết quả.

Trong phân tích độ phức tạp, cần chứng minh tổng độ dài cạnh của mọi khối là
$O(n\log n)$. Sự thật này có thể chứng minh dưới điều kiện hàm biên là đơn
điệu; biên của ellipse có thể tách thành bốn hàm đơn điệu. Ý tưởng chứng
minh là: với mỗi tọa độ $x$, các khối theo tọa độ $y$ tương ứng chỉ gồm một
số hằng các khối cộng thêm và một số khối khác; với cùng một kích thước
khối, tổng độ dài của các khối được cộng thêm không thể vượt quá $n$. Vì có
$O(\log n)$ loại độ dài, kết luận suy ra.

\section{1007 Power Station of Art}

Trước hết, thao tác là khả nghịch. Điều này có nghĩa là từ đầu đến cuối ta
chỉ cần thao tác trên đồ thị thứ nhất. Tiếp theo, hai đồ thị có thể trở nên
bằng nhau khi và chỉ khi từng thành phần liên thông đều có thể trở nên bằng
nhau.

Xét một thành phần liên thông. Trước hết, multiset các số trong thành phần
này ở hai đồ thị phải giống nhau. Sau đó nhận thấy cách hoán đổi trong đề có
thể hiểu tương đương như sau: hoán đổi cả số và màu của hai đỉnh, rồi lật cả
hai màu. Kết quả thao tác như vậy giống hệt thao tác ban đầu. Trong cách hiểu
này, nếu một số bị hoán đổi số lần lẻ thì màu của nó bị lật; nếu bị hoán đổi
số lần chẵn thì màu không đổi. Ta chia trường hợp.

\begin{itemize}
  \item Nếu đồ thị là đồ thị hai phía, tô đen trắng đồ thị đó. Với một số
  bất kỳ, nếu biết màu đen trắng ban đầu của nó và màu của đỉnh, thì sau khi
  chuyển số đó đến bất kỳ vị trí nào, màu tương ứng của đỉnh là xác định.
  Hơn nữa, vì đồ thị liên thông, các số có thể được hoán vị tùy ý đến mọi
  đỉnh. Khi đó điều kiện là: chia các số theo giá trị
  ``màu đỉnh xor màu số'' thành hai multiset, và hai đồ thị phải có hai
  multiset tương ứng giống nhau.
  \item Nếu đồ thị không phải đồ thị hai phía, thành phần liên thông có một
  chu trình lẻ. Khi đó điều kiện là multiset mọi số trong hai đồ thị phải
  giống nhau, và tính chẵn lẻ của số lượng màu không thể thay đổi.

  Lý do là có thể xem đồ thị như một đồ thị hai phía cộng thêm một số cạnh
  thừa. Với bất kỳ chu trình lẻ nào, ta dựng được một phương án giữ nguyên
  mọi số và chỉ lật màu ở hai vị trí, bất kể hai vị trí đó ban đầu có màu
  gì. Cách dựng: đánh số các đỉnh trên chu trình lẻ từ $1$ đến $n$; trước hết
  thực hiện $n-1$ thao tác để chuyển số trên đỉnh $1$ đến đỉnh $n$, sau đó
  thực hiện một thao tác hoán đổi số trên đỉnh $1$ và đỉnh $n$, rồi dùng
  $n-2$ thao tác để chuyển số trên đỉnh $n$ đến đỉnh $2$.
\end{itemize}

\section{1008 Command and Conquer: Red Alert 2}

Xét một tay bắn tỉa có tầm $k$; phạm vi tấn công của người đó là một khối lập
phương. Ta có thể phát biểu lại như sau: từ vị trí $(x,y,z)$, tay bắn tỉa có
thể tấn công mọi điểm $(x',y',z')$ thỏa
\[
  x\le x'\le x+2k,\qquad
  y\le y'\le y+2k,\qquad
  z\le z'\le z+2k.
\]

Khi đó, nếu tọa độ của tay bắn tỉa ở một chiều nào đó nhỏ hơn tọa độ nhỏ nhất
theo chiều đó trong các điểm còn lại, ta nên di chuyển tay bắn tỉa đến đúng
tọa độ nhỏ nhất đó, vì điều này chỉ làm tăng số mục tiêu có thể bắn. Khi cả
ba tọa độ đều bằng tọa độ nhỏ nhất tương ứng, tay bắn tỉa bắt buộc phải tiêu
diệt toàn bộ quân địch có tọa độ nhỏ nhất ở một trong ba chiều thì mới có thể
tiếp tục di chuyển. Lúc này tham lam chọn tiêu diệt nhóm quân địch yêu cầu
$k$ nhỏ nhất.

Độ phức tạp thời gian là $O(n\log n)$.

\section{1009 Typing Contest}

Để tránh số thực, trước hết nhân $f$ với $100$. Khi đó công thức tính tốc độ
của học sinh $1$ là
\[
  \frac{s_1}{10000}\left(10000-f_1f_2-f_1f_3-\cdots-f_1f_k\right).
\]
Để trình bày gọn hơn, phần sau bỏ qua hệ số $\frac{1}{10000}$ phía trước.

Công thức tốc độ của học sinh $1$ có thể viết lại thành
\[
  s_1\left(10000-f_1\left(\sum_{i\in S} f_i-f_1\right)\right),
\]
trong đó $S$ là tập học sinh được chọn. Nếu cố định được
$\sum_{i\in S} f_i$, thì tốc độ của từng học sinh cũng được cố định. Vì vậy
ta liệt kê
\[
  \sum_{i\in S} f_i=F.
\]
Bài toán trở thành bài ba lô $01$: mỗi học sinh có trọng lượng $f_i$, giá trị
\[
  s_i\left(10000-f_i(F-f_i)\right),
\]
và tổng dung lượng của ba lô là $F$.

Nếu liệt kê $F$ từ $1$ đến $\sum_{i=1}^{n} f_i$ thì không chấp nhận được.
Thực ra chỉ cần liệt kê đến $100\sqrt n+2$; chứng minh ở dưới. Vì vậy độ
phức tạp thời gian là $O(5000n^2)$.

Chứng minh: giả sử chọn $k$ vật. Với một vật bất kỳ $t$, ta có
\[
  s_t\left(10000-f_t\left(\sum_{i=1}^{k}f_i-f_t\right)\right)>0,
\]
ở đây $f_i$ là giá trị ban đầu sau khi đã nhân $100$. Do đó
\[
  f_t\sum_{i=1}^{k}f_i-f_t^2<10000.
\]
Lấy tổng theo $t=1,2,\ldots,k$ ta được công thức A:
\[
  \left(\sum_{i=1}^{k}f_i\right)^2-\sum_{i=1}^{k}f_i^2<10000k.
\]

Từ
\[
  f_t\sum_{i=1}^{k}f_i-f_t^2<10000
\]
suy ra
\[
  f_t^2\sum_{i=1}^{k}f_i-f_t^3<10000f_t.
\]
Lấy tổng theo $t=1,2,\ldots,k$ ta được công thức B:
\[
  \sum_{i=1}^{k} f_i^2
  < 10000+\frac{\sum_{i=1}^{k}f_i^3}{\sum_{i=1}^{k}f_i}
  \le 20000.
\]
Kết hợp A và B:
\[
  \left(\sum_{i=1}^{k}f_i\right)^2
  < \sum_{i=1}^{k}f_i^2+10000k
  < 10000k+20000.
\]
Vì vậy
\[
  f_1+f_2+\cdots+f_k
  < 100\sqrt{k}+2
  \le 100\sqrt n+2.
\]

\section{1010 Array}

Nên vẽ một bảng ô vuông $n\times n$, trong đó ô $(i,j)$ biểu diễn đoạn con
$A[i\ldots j]$. Với mọi cặp gần nhất có cùng giá trị $a_i=a_j$, tô đen hình
vuông có góc trái dưới là $(i+1,i+1)$ và góc phải trên là $(j-1,j-1)$.

Đặt
\[
  u_i=\max\{j\mid j<i,\ a_j=a_i\},\qquad
  v_i=\min\{j\mid j>i,\ a_j=a_i\}.
\]
Đặc biệt, nếu trước $a_i$ không có số bằng nó thì $u_i=0$; nếu sau $a_i$
không có số bằng nó thì $v_i=n+1$.

Bây giờ xét $U,V$ tương ứng với dãy $B$. Chắc chắn có $u_{b_1}=0$. Với mọi
$i\in [1,n-1]$, nếu $b_i<b_{i+1}$ thì
\[
  v_i=b_{i+1},\qquad u_{b_{i+1}}=i,
\]
tức $a_i=a_{b_{i+1}}$. Riêng khi $b_{i+1}=n+1$ thì chỉ có $v_i=n+1$, vì
$u_{n+1}$ không được định nghĩa.

Ngoài các vị trí trong $U,V$ đã bị $B$ chỉ định, những vị trí khác có thể
điền tùy ý. Tuy nhiên, vì với $u_i\ne 0$ thì $a_i=a_{u_i}$, và với
$v_i\ne n+1$ thì $a_i=a_{v_i}$, nên với mọi $i\in[1,n]$ cần có
\[
  u_i<i,\qquad v_i>i,\qquad u_{v_i}=i,\qquad v_{u_i}=i,
\]
đồng thời $b_i>v_i$ trong ràng buộc tương ứng. Các điều kiện này cũng chính
là những gì biên bậc thang của dãy $B$ thể hiện trên hình ô vuông.

Sau đó có thể nhận ra nửa tam giác trên của bảng ô vuông là một đường gấp
khúc dạng bậc thang, và dãy $B$ tương ứng một-một với đường gấp khúc này.
Các vị trí bị chỉ định trong $u,v$ chính là phần biên ngang và biên dọc của
đường gấp khúc; cách nhìn này giúp trực quan hóa các quy tắc trên.

Khi đã hiểu phần trên, bài toán chuyển thành: với mỗi vị trí từ $1$ đến $n$,
có một ổ cắm trái, tức mảng $U$, và một ổ cắm phải, tức mảng $V$. Ta cần dùng
một số dây để ghép cặp các ổ cắm trái-phải, biểu diễn quan hệ bằng nhau. Nếu
một dây nối ổ cắm trái của $i$ với ổ cắm phải của $j$, ký hiệu dây này là
$T(i,j)$, nghĩa là $a_i=a_j$.

Dây $T(i,j)$ phải thỏa $i<j$. Vị trí $0$ có vô hạn ổ cắm phải; vị trí $n+1$
có vô hạn ổ cắm trái. Các vị trí khác mỗi nơi có đúng một ổ cắm trái và một
ổ cắm phải. Mọi ổ cắm trên các vị trí $1$ đến $n$ phải được cắm đầy; các ổ
cắm ở $0$ và $n+1$ thì không bắt buộc. Một số dây đã bị chỉ định sẵn. Ngoài
ra, $T(i,j)$ phải thỏa
\[
  c_j<i<j,
\]
trong đó mảng $C$ có thể tính từ dãy $B$; định nghĩa của $C$ khá rõ trên
hình ô vuông.

Duy trì hai tập ổ cắm phải $A,B$: $A$ là tập các ổ cắm bắt buộc phải được
cắm, còn $B$ là tập các ổ cắm cắm hay không đều được. Ban đầu $B$ chứa vô hạn
ổ cắm ở vị trí $0$. Duyệt $i$ từ $1$ đến $n$. Nếu ổ cắm trái của $i$ chưa bị
chỉ định, ta xử lý như sau:
\begin{itemize}
  \item Nếu $A$ rỗng, tìm trong $B$ ổ cắm phải có chỉ số nhỏ nhất và thỏa
  điều kiện của mảng $C$ để ghép cặp với nó. Nếu không tìm được thì in
  \texttt{NO}.
  \item Ngược lại, nếu trong $B$ không có ổ cắm phải có thể ghép với nó, trực
  tiếp tìm trong $A$ ổ cắm phải phù hợp có chỉ số nhỏ nhất để ghép.
  \item Nếu cả $A$ và $B$ đều có ứng viên, gọi $x,y$ lần lượt là chỉ số nhỏ
  nhất của ổ cắm phải phù hợp trong $A,B$. Xóa $x,y$ khỏi hai tập tương ứng
  và đưa $\max(x,y)$ vào $B$.
\end{itemize}

Sau đó, nếu ổ cắm phải của $i$ chưa bị chỉ định, thêm nó vào tập $A$. Cuối
cùng, vì còn vô hạn ổ cắm trái ở vị trí $n+1$, chỉ cần kiểm tra các phần tử
còn lại trong $A$ có đều lớn hơn $c_{n+1}$ hay không. Nếu có thì tất cả đều
có thể ghép với $n+1$, đáp án là \texttt{YES}. Trong quá trình trên, hễ không
tìm được ứng viên cần thiết thì đáp án là \texttt{NO}.

Độ phức tạp thời gian là $O(n\log n)$. Có thể làm $O(n)$, nhưng không cần.

\section{1011 Game}

Nếu từ một điểm $(i,j)$, mọi điểm có thể đến bằng đúng một bước hoặc đúng hai
bước đều không phải điểm đặc biệt, thì sau một phân loại đơn giản, trạng thái
thắng thua của $(i+1,j-1)$ giống với trạng thái của $(i,j)$.

Số điểm đặc biệt chỉ là $O(n)$, nên số điểm không thể quy giản về
$(i+1,j-1)$ cũng chỉ là $O(n)$. Lấy toàn bộ các điểm này ra rồi chạy tìm kiếm
có nhớ. Với điểm thường, đi theo hướng giảm $i+j$ để tìm điểm đầu tiên không
thể quy giản; tại điểm đó thì duyệt thô hai loại chuyển trạng thái. Độ phức
tạp là
\[
  O((n+q)\log(n+q)).
\]

Tất nhiên, hằng số của tìm kiếm có nhớ hơi lớn. Dùng quét đường thẳng có thể
giảm hằng số khoảng năm lần, nhưng vì lời giải chuẩn dùng tìm kiếm có nhớ nên
giới hạn thời gian được đặt rộng hơn.

\end{document}
